\section{The Mixture-of-Models Artifact and Lifecycle}
\label{sec:artifact}

A model architecture needs a model artifact, not merely runtime configuration.
For MoM, that artifact is the \emph{bundle}: a typed envelope that gives a
composite model one interface, identity, policy, and finite execution semantics
while leaving constituent checkpoints in their native formats. The same logical
object must remain identifiable across authoring, optimization, evaluation, and
deployment. The bundle carries the logical half of the co-design; binding and
locking instantiate its physical half, and realized traces provide common
evidence for optimizing both. We therefore define six lifecycle stages, separating policy and
learned assets from environment-specific endpoints and traces:
\begin{equation}
 \text{source}\rightarrow\text{typed IR}\rightarrow\text{bundle}
 \rightarrow\text{binding}\rightarrow\text{lock}
 \rightarrow\text{realized trace}.
 \label{eq:artifact-pipeline}
\end{equation}

\begin{figure}[H]
  \centering
  \resizebox{\linewidth}{!}{\begin{tikzpicture}[
  x=1cm,y=1cm,font=\sffamily,>=Latex,
  flow/.style={-{Latex[length=1.7mm]},draw=MoMGray,line width=0.65pt},
  group/.style={draw=MoMBorder,fill=white,rounded corners=3pt,line width=0.45pt},
  portable/.style={group,fill=MoMBlueLight!20},
  base/.style={draw=MoMBorder,fill=white,rounded corners=2.5pt,
    line width=0.58pt,minimum height=1.02cm,inner xsep=6pt,
    inner ysep=3pt,align=center,text=MoMText},
  source/.style={base,draw=none,fill=MoMBlueLight},
  bundle/.style={base,draw=none,fill=MoMDark,text=white},
  trace/.style={base,draw=MoMDark,fill=MoMBlue,line width=0.72pt},
  header/.style={font=\sffamily\bfseries\fontsize{7.0}{8.0}\selectfont,
    text=MoMMuted}
]

% Six lifecycle stages share one baseline; the semantic boundary lies between
% the portable bundle and the environment-owned binding.
\draw[portable] (0.05,0.15) rectangle (7.25,2.62);
\draw[group] (7.55,0.15) rectangle (13.70,2.62);
\node[header,anchor=west] at (0.36,2.30) {Portable artifact representations};
\node[header,anchor=west] at (7.86,2.30) {Environment-specific realization};

\node[source,minimum width=1.70cm] (src) at (1.05,1.22)
  {{\bfseries\fontsize{8.0}{8.8}\selectfont Source}\\[-1pt]
   {\fontsize{6.5}{7.2}\selectfont\color{MoMMuted} DSL / YAML}};

\node[base,minimum width=1.70cm] (ir) at (3.35,1.22)
  {{\bfseries\fontsize{8.0}{8.8}\selectfont Typed IR}\\[-1pt]
   {\fontsize{6.5}{7.2}\selectfont\color{MoMMuted} semantics}};

\node[bundle,minimum width=1.95cm] (bundle) at (5.83,1.22)
  {{\bfseries\fontsize{8.0}{8.8}\selectfont MoM bundle}\\[-1pt]
   {\fontsize{6.5}{7.2}\selectfont model artifact}};

\node[base,minimum width=1.55cm] (binding) at (8.35,1.22)
  {{\bfseries\fontsize{8.0}{8.8}\selectfont Binding}\\[-1pt]
   {\fontsize{6.2}{6.9}\selectfont\color{MoMMuted} fleet policy}};

\node[base,minimum width=1.55cm] (lock) at (10.25,1.22)
  {{\bfseries\fontsize{8.0}{8.8}\selectfont Lock}\\[-1pt]
   {\fontsize{6.2}{6.9}\selectfont\color{MoMMuted} revisions}};

\node[trace,minimum width=1.90cm] (realized) at (12.50,1.22)
  {{\bfseries\fontsize{8.0}{8.8}\selectfont Realized trace}\\[-1pt]
   {\fontsize{6.3}{7.0}\selectfont\color{MoMText} calls / outcomes}};

\draw[flow] (src) -- (ir);
\draw[flow] (ir) -- (bundle);
\draw[flow] (bundle) -- (binding);
\draw[flow] (binding) -- (lock);
\draw[flow] (lock) -- (realized);

\end{tikzpicture}
}
  \caption{\textbf{The MoM lifecycle separates portable representation
  from environment-specific realization.} The lifecycle engine compiles source into
  typed IR and an immutable bundle, then binds and resolves logical references
  into a lock. Every realized trace is attributed to the
  bundle, binding, and lock; credentials and live health remain outside the
  portable artifact.}
  \label{fig:artifact-lifecycle}
\end{figure}

\subsection{Six stages, six responsibilities}

The following definitions are normative requirements of the proposed format,
not a description of the current \vsr{} artifact format.

\paragraph{Authoring source.}
Source is a human-readable DSL or YAML with comments, fragments, and
defaults.  Sources that normalize to the same IR define the same logical model.

\paragraph{Canonical typed IR.}
The IR is the semantic authority. It declares schemas and entry points, logical
model requirements, signals and projections, action graphs, extension
interfaces, contracts, behavior variants, request-preference domains, and
bounds. It contains neither credentials
nor implicit defaults. Canonical maps, numbers, and asset references enable
stable hashes and semantic diffs, and unknown fields are rejected. Model
requirements have no competing manifest authority; hard constraints are typed
fail-closed predicates; and every contract check names explicit success and
failure successors.

\paragraph{Portable bundle.}
The bundle is an immutable distribution unit around the IR:
\begin{verbatim}
manifest.json             interface, content index, compatibility
mom.ir.json               canonical executable specification
source/ assets/           authoring sources and semantic assets
contracts/ eval/          schemas and evaluation protocols
targets/ MODEL_CARD.md    realization templates and documentation
\end{verbatim}
Source is optional, allowing an author to export the executable object without
revealing its authoring representation. Leaf artifacts may be safetensors,
ONNX, GGUF, or immutable repository objects; closed constituents remain typed
provider references. Large weights may be external but content-addressed.
Plugin code is not implicitly trusted: the bundle declares an operator set,
interface, effects, and digest; deployment policy allows or rejects it.

A \emph{behavior variant} such as accuracy-first or security-first belongs to
logical semantics and therefore to model identity; it declares the default and
allowed domain of request preferences. A \emph{realization profile}
such as CUDA, ROCm, CPU, edge, or hybrid instead declares portable runtime and
accelerator requirements. It may be distributed with the bundle, but it contains
no endpoint, secret, device identifier, or live capacity and does not change
logical identity.

Backend compatibility is not a scalar \texttt{backend} choice. A binding keeps
six axes distinct: wire transport, API adapter, model runtime, leaf-artifact
format, accelerator, and deployment driver. Thus one admissible trace may use a
ROCm-hosted vLLM model, a CUDA specialist, a CPU verifier, and a managed cloud
API without assigning the composite model itself a single hardware backend.

\paragraph{Deployment binding.}
A binding is an environment-owned mapping
\begin{equation}
 \mathcal{B}:\mathcal{V}\cup\mathcal{Q}\rightarrow 2^{\mathcal{E}}
 \end{equation}
from logical model or operator references to candidate endpoints or local
deployments under a selected realization profile. For
$r\in\mathcal{V}\cup\mathcal{Q}$, each
$d\in \mathcal{B}(r)$ must
satisfy modality, context, capability, data-policy, and interface requirements;
operator candidates also satisfy their declared effect and sandbox
requirements. Bindings declare placement, replicas, limits, prices, and
fallbacks, but contain only secret references. Each candidate exposes its
capability evidence, capacity envelope, and pricing assumptions, so resolution
cannot infer them silently from an endpoint name. Here $\mathcal{E}$ is the
environment's set of candidate deployments.

\paragraph{Resolution lock.}
Resolution lock $L$ records the complete resolved candidate inventory: model,
tokenizer, adapter, image, runtime, accelerator, protocol-adapter, asset, and
plugin identities plus capability evidence. Binding says what may run; lock says
what was resolved; the trace says what actually ran. Revision evidence is typed
as \emph{pinned} when an immutable content or repository identifier can be
recorded and \emph{observed-opaque} when a provider exposes only an alias. The
latter retains an observation time and a content-addressed probe record without
claiming bitwise reproducibility. Floating references are a development mode,
not a production lock.

\paragraph{Realized trace.}
The trace records semantic/bundle/binding/lock digests, behavior variant,
request preference, realization profile, constraint
checks, signals, actions, resolved calls, timing, tokens, usage, accounting
version, realized cost, retries, contracts, and terminal status. Content may be
encrypted, redacted, hashed, or referenced under data policy. A trace is an
observation, not model definition.

\subsection{Identity and portability}

A human coordinate such as \texttt{modelId:version} is not an integrity
identity. We instead separate four content identities. The
\emph{semantic digest} commits to canonical IR and its path-sorted semantic
assets. The \emph{bundle digest} additionally commits to the manifest, source,
contracts, evaluation protocols, target templates, and documentation. The
\emph{binding digest} identifies canonical environment policy, while the
\emph{lock digest} identifies one completed resolution. The manifest stores the
semantic digest; the external lock stores
bundle and binding digests; the lock digest is computed only after the lock is
complete. No object therefore contains its own hash.

Bindings, credentials, and traces are excluded from logical identity.
Rebinding realizes the same logical model but produces new binding and lock
identities; changing a controller, prompt, contract, requirement, behavior
variant, or topology changes the semantic digest. Changing only source
formatting, a model card, or a realization template preserves logical identity
but changes the distributed bundle. Every result carries all four digests.
Together they distinguish semantic,
distribution, environment, and realized-resolution equivalence.

Portability preserves logical and control semantics, not predictive equivalence
across substitutions. It is tested by fixing semantic and bundle digests,
resolving them in two environments, and reporting quality, validity, cost, and
latency under both binding and lock identities. Capability substitutions are
explicit and disclosed; two edited recipes are not portable merely because
they share an API path.

\subsection{Build, import, validation, and export}

Build expands source defaults into canonical IR and performs four classes of
checks: schema and type validity; reference, reachability, and bound
correctness; well-formed contract branches and statically decidable constraint
consistency; and path, digest, and license-record integrity. Explicit migration
handles older source versions; runtime evidence resolves guards that cannot be
decided statically. Import instead verifies an already built artifact, records
its content identity, and never recompiles an optional source representation.
At deployment, every reachable call needs an admissible binding or a declared
fallback or abstention path.

Export writes canonical IR, a deterministic manifest, content-addressed assets,
and evaluation protocols---never credentials, endpoint health, online counters,
or private traces. It may produce a thin artifact of pinned references, a
materialized artifact containing redistributable leaf objects, or a targeted
artifact with selected runtime objects. Materialization fails rather than
pretending that a closed provider is portable. Source round-tripping preserves
semantics, not formatting. IR diff reports behavior changes such as a widened
candidate set or relaxed residency guard.

Consumers verify digests, schemas, operator sets, and engine capabilities;
deployment policy
enforces plugin and license rules and probes model access. Evaluation results,
build provenance, SBOMs, and signatures are digest-bound attestations outside
the bundle, so new evidence does not redefine the evaluated model.

This boundary follows useful precedents without reducing MoM to any one of
them. Hugging Face repositories motivate a save/load model coordinate and
immutable revision pinning; MLflow motivates signatures shared by loading,
serving, and evaluation; ONNX separates a versioned logical IR from
capability-negotiated execution providers; and OCI supplies content-addressed
descriptors and distribution primitives
\citep{huggingface_2026_models,mlflow_2026_model_format,
onnx_2026_ir,onnxruntime_2026_execution_providers,oci_2025_image_spec}.
MoM applies this separation at the model-call level: an accelerator or runtime
realizes the model but does not define it.

\subsection{One lifecycle across training, evaluation, and inference}

An optimizer consumes data or traces from a deployed realization
$(\mathcal{M}_k,\mathcal{B}_k,L_k)$ and emits a complete logical or binding
candidate, not an opaque mutation. Under recorded resolution locks, evaluation
compares candidate and incumbent on quality, violations, cost, latency, retries,
and uncertainty before promotion. Serving emits traces, monitoring detects
drift, replay reconstructs decisions, and approved traces feed offline
optimization. Immutable local revisions support revision-pinned re-execution;
stochastic sampling and runtime scheduling still require recorded seeds and
events. Opaque provider aliases support attributed re-evaluation and drift
measurement, not bitwise reproduction.

Online adaptation may update trace-linked local state, but the model changes
only when that state is promoted into a hashed asset or IR revision.  Session
protection and fallback remain runtime state, preventing one replica's traffic
from silently redefining another's artifact.
